\documentclass[aps,twocolumn, superscriptaddress, english, nofootinbib, 10pt]{revtex4-2}

\usepackage{float}

\usepackage{amssymb, amsmath}
\usepackage[utf8]{inputenc}
\usepackage{subfigure}
\usepackage{comment}
\usepackage{graphicx}
\usepackage{hyperref}
\usepackage{enumerate}
\usepackage{booktabs}
\usepackage{lipsum}
\usepackage{xspace}
\usepackage{aas_macros}
\usepackage{multirow}

\usepackage[usenames,dvipsnames]{xcolor}
\hypersetup{
    colorlinks = true,
    citecolor = {MidnightBlue},
    linkcolor = {BrickRed},
    urlcolor = {BrickRed}
}


\newcommand{\nv}{\hat{\bf n}}
\newcommand{\wtj}[6]{\left(\begin{array}{ccc} #1 & #2 & #3\\#4 & #5 & #6\end{array} \right)}
\newcommand{\nmt}{{\tt NaMaster}\xspace}
\newcommand{\todo}[1]{{\color{blue}\bf TODO: #1}}


\begin{document}


\title{Analytical covariances for catalogue-based pseudo-$C_\ell$s}

\author{Kevin Wolz}
\email{kevin.wolz@physics.ox.ac.uk}
\affiliation{Department of Physics, University of Oxford, Denys Wilkinson Building, Keble Road, Oxford OX1 3RH, United Kingdom}
\author{David Alonso}
\affiliation{Department of Physics, University of Oxford, Denys Wilkinson Building, Keble Road, Oxford OX1 3RH, United Kingdom}
\author{Elyas Farah}
\affiliation{Argelander-Institut für Astronomie, Universität Bonn, Auf dem Hügel 71, D-53121 Bonn, Germany}
\author{Robert Reischke}
\affiliation{Argelander-Institut für Astronomie, Universität Bonn, Auf dem Hügel 71, D-53121 Bonn, Germany}


\author{Add yourselves!}
\affiliation{And your affiliations!}

\begin{abstract}
  Multiple cosmological observables, such as the galaxy overdensity or cosmic shear, consist of fields sampled at the discrete positions of astrophysical sources. In recent work, we presented a method to estimate the angular power spectra of such fields, avoiding the construction of pixelised sky maps and the finite-resolution effects associated with them. In this work, we present a method to estimate the disconnected (also known as ``Gaussian'') covariance of these angular power spectra, addressing subtle effects, including the effective area overlap between different catalogue-based fields and the additional Poisson-like variance arising from the discrete nature of the catalogues. The method relies on the so-called Narrow-Kernel Approximation to account for the contribution of distinct source pairs to the estimator, while including the noise-like contributions from self-pairs exactly. We explicitly compare this approach with a brute-force method that can produce the exact covariance for sparse samples. We show that the method is accurate in realistic scenarios, spanning both dense and noise-dominated datasets (e.g., cosmic shear) and sparse, signal-dominated observables (e.g., fast radio bursts). The method is implemented in the public code \nmt.
\end{abstract}

%  \date{\today}
\maketitle

\section{Introduction}\label{sec:intro}
    Projected fields of cosmological observables contain a wealth of information about the anisotropies in the late-time Universe, the growth of structure, the distribution of circumgalactic gas in the Universe. Example tracers include galaxy clustering, gravitational lensing of galaxy shapes, galaxy velocities, kinematic and thermal Sunyaev-Zel'dovich effects, fast radio burst dispersion, etc. There is particular interest in cross-correlating tracers to break cosmological and astrophysical degeneracies and be insensitive to additive systematic bias. Analysing these observables across the sky requires robust estimators with known theoretical predictions and statistical distributions, at least approximately. An example is the pseudo-$C_\ell$ (PCL) estimator of the angular power spectrum, a fast, powerful, and therefore widely used tool to analyse auto- and cross-correlations on the partial sky. This estimator, originally applied to pixel-level maps, was recently extended to catalogue-level fields, thereby avoiding biases from aliasing and mode-coupling instability [cites]. This work extends the work in \todo{cite} to covariances for catalogue-level PCLs, using the improved narrow-kernel approximation (iNKA) introduced in [cite]. We also present a brute-force method that evaluates catalogue-based covariances without approximations.

    Catalogue-level PCL estimation is the natural way to deal with discrete tracers. In the traditional way, one would bin sources into pixels, which has three main disadvantages:
    \begin{itemize}
        \item[a)] it can introduce nonanalytical pixel window functions requiring expensive simulations \todo{cites},
        \item[b)] it causes shot-noise bias to the field estimate,
        \item[c)] it causes shot-noise bias to the mask estimate, introducing instability in the mode decoupled estimate.
    \end{itemize}
    Those effects can be avoided by evaluating the spherical harmonic transforms directly at the source positions using \todo{method name and cite}. As presented in \todo{cite}, problems b) and c) are eliminated before mode decoupling by subtracting, an analytical shot-noise estimate at the field and mask level, respectively, which has the practical side effect of making the estimator immune to any white-noise-like bias.
    
    This algorithm has been implemented in the public code NaMaster[cites]. The code supports different field types for catalog-level observables:
    \begin{itemize}
        \item tracer fields (\texttt{NmtFieldCatalog}). These are fields that are associated with an underlying (in principle continuous, in principle general-spin) field whose values are sampled at the discrete tracer positions, and whose mask is defined by those positions. Examples are: galaxy shear, FRB dispersion measure, etc.
        \item momentum fields (\texttt{NmtFieldCatalogMomentum}). These fields, recently introduced in \todo{cite}, are associated with a number density contrast, evaluated through source number counts per sky area, and weighted by a local tracer field evaluated at the source positions. The mask can be provided either through a random catalog, or though a pixelated completeness mask. An example is the galaxy momentum density (galaxy number density weighted by radial velocity) used for kinematic Sunyaev-Zel'dovich power spectra introduced in \todo{cite} and used in \todo{cite}.
        \item galaxy clustering fields (\texttt{NmtFieldCatalogClustering}). A special case of the previous case, with unit tracer weights at each source position. These are fields that are associated with a number density contrast, evaluated through source number counts per sky area. The mask can be provided either through a random catalog, or through a pixelated completeness mask. An example is galaxy clustering.
    \end{itemize}
    Another recent work \todo{cite} introduced the option of linearly deprojecting systematics templates at the catalog level. With this covariance algorithm, iNKA cross-tracer covariances between any combination of catalog- and map-based fields (\texttt{NmtField)} are supported. We validate the code and demonstrate its application on realistic mock data the algorithm in \todo{Section}. 

    Most applications of p$C_\ell$s to astrophysical inference assume a Gaussian likelihood. Therefore, knowing the covariance is key to performing robust inference with p$Cl\ell$s. \todo{summarize covariances literature, also mention super-sample covariance}. This work builds on the iNKA approximation for disconnected Gaussian fields, which assumes that the harmonic-space two-point correlators of two masked fields $a$ and $c$ satisfy the narrow-kernel approximation, i.e., they can be factorized as
    \begin{equation*}
        \langle \tilde{a}_{\ell m} \tilde{c}_{\ell' m'} \rangle \approx (C_\ell^{ac} + C_{\ell'}^{ac})\mathcal{M}_{\ell m,\ell' m'}(w_aw_c)/2
    \end{equation*}
    In addition, pulling out $C_\ell$ out of the integral ignores the mask-induced mode coupling, which can be accounted for by replacing $C_\ell$ by its normalized mode-coupled version (improved NKA). The iNKA has the advantage of being an $\mathcal(\ell_{\rm max}^3)$ operation, which guarantees computations stay permissive up to $\ell_{\rm max}\approx 10,000$, relevant for surveys targeting small-scale information, such as galaxy clustering and cosmic shear. In this work, we also present a JAX-based implementation of an exact ``brute-force'' method, relevant for low-number density tracers where small scales are dominated by shot noise. We demonstrate an application of this brute-force method to FRB-like simulations with \todo{maximum multipole}.

    In this work, we promote the formalism for the pseudo-$C_\ell$ covariance to the case of discrete sources. In the iNKA logic, a key challenge is the multiplication of catalog masks entering the normalization and the mode-coulpling kernel. In the case of catalog auto-spectra, these mask products introduce shot noise that add to the total covariance, which we identify with additional ``noise-noise'' and ``noise-signal'' components. These are disconnected trispectra that result from inhomogeneously sampling the sky. In the case of cross-spectra, mask products, appearing in the signal-signal part, are formally zero due to non-overlapping sources. We circumvent this issue by smoothing the catalog masks in harmonic space at a scales above the mean inter-particle distance of the catalog. In practice, this means we assume that the catalogs are inhomogeneously sample Gaussian fields, ignoring the non-Gaussian distribution of the underlying point process. \todo{elaborate/mention this in the conclusions as a potential are to improve in future work} We find this approximation holds well for all our validation cases. \todo{Mention where this might break down (sparse catalogs, such as FRBs etc.)}.

    This paper is structured as follows: \todo{add}.

\section{Crash course on pseudo-$C_\ell$s}\label{sec:pcls}
  We begin with a short introduction to the pseudo-$C_\ell$ estimator to set the scene, establish the notation that will be used in the rest of the paper, and clarify the challenges encountered in extending the estimation of power spectrum covariances to the case of discretely-sampled fields. For simplicity, we focus our discussion in the main text on scalar fields, and we summarise the modifications to the method presented here in the case of spin-$s$ fields in Appendix \ref{app:spin_s}. We refer readers to previous papers presenting this formalism for further details. In particular \todo{cite} presents the general pseudo-$C_\ell$ estimator, \todo{cite} and \todo{cite} present the improved Narrow Kernel Approximation (iNKA) used to estimate pseudo-$C_\ell$ covariances, and \todo{cite} introduces catalogue-based power spectra.

  \subsection{Pseudo-$C_\ell$s}\label{ssec:pcls.pcl}
    Consider a field $a(\nv)$ defined on the sphere, where the unit vector $\nv$ denotes a given position on the celestial sphere. In all practical cases, we observe a ``masked'' version of this field $\tilde{a}\equiv\,w_a\,a$, where $w_a(\nv)$ is the sky mask associated with $a$. In the simplest case, $w_a$ is a binary map that simply defines the regions of the sky where the field has been observed. More generally, $w_a$ can be thought of as a weights map that optimally combines the different sky regions to maximise the signal-to-noise ratio of the field's power spectrum. The spherical harmonic transform of $\tilde{a}$ and $a$ are related to each other through a convolution:
    \begin{equation}\label{eq:mask_harmonic}
      \tilde{a}_{\ell m}=\sum_{\ell'm'}{\cal M}_{\ell m,\ell'm'}(w_a)\,a_{\ell'm'},
    \end{equation}
    where the map-level coupling coefficients ${\cal M}$, depending on the sky mask, are
    \begin{equation}\label{eq:map_mcm}
      {\cal M}_{\ell m,\ell'm'}(w_a)\equiv\int d\nv\,w_a(\nv)\,Y^*_{\ell m}(\nv)\,Y_{\ell'm'}(\nv).
    \end{equation}

    In the case of full-sky observations ($w_a=1$ throughout, and ${\cal M}_{\ell m,\ell'm'}(w_a)=\delta^K_{\ell\ell'}\delta^K_{mm'}$, where $\delta^K_{ij}$ is the Kronecker delta), the optimal estimator for the power spectrum of two fields, $a$ and $b$, is
    \begin{equation}
      \hat{C}^{ab}_\ell=\frac{1}{2\ell+1}\sum_m{\rm Re}(a^*_{\ell m}b_{\ell m}).
    \end{equation}
    The pseudo-$C_\ell$ estimator consists of na\"ively applying this to the masked fields instead:
    \begin{equation}\label{eq:pcl_def}
      \tilde{C}^{ab}_\ell\equiv\frac{1}{2\ell+1}\sum_m{\rm Re}(\tilde{a}^*_{\ell m}\tilde{b}_{\ell m}).
    \end{equation}
    However, just like the presence of a mask induces a coupling between different angular multipoles at the map level (Eq. \ref{eq:mask_harmonic}), so does it cause mode-coupling at the level of the power spectrum. Indeed, the ensemble average of the pseudo-$C_\ell$ estimator is:
    \begin{equation}\label{eq:pcl}
      \left\langle \tilde{C}^{ab}_\ell\right\rangle=\sum_{\ell'} M_{\ell\ell'}(w_a,w_b)\,C^{ab}_{\ell'},
    \end{equation}
    where $C^{ab}_\ell$ is the true underlying power spectrum, and $M_{\ell\ell'}(w_a,w_b)$ is the \emph{mode-coupling matrix}. The latter depends only on the pseudo-$C_\ell$ of the two masks, $\tilde{C}^{w_a,w_b}_\ell$, and can be calculated analytically with modest computational complexity ($\sim{\cal O}(\ell_{\rm max}^3)$). Specifically, for spin-0 quantities, $M_{\ell\ell'}(w_a,w_b)=(2\ell'+1)\Xi_{\ell\ell'}(w_a,w_b)$, where the coupling coefficients $\Xi_{\ell\ell'}$ are
    \begin{equation}\label{eq:Xi}
      \Xi_{\ell\ell'}(w_a,w_b)\equiv\sum_{\ell''}\frac{2\ell''+1}{4\pi}\,\tilde{C}^{w_a,w_b}_{\ell''}\,\wtj{\ell}{\ell'}{\ell''}{0}{0}{0}^2,
    \end{equation}
    involving the Wigner-$3j$ symbols.

  \subsection{Gaussian covariances}\label{ssec:pcls.cov}
    The statistical uncertainties of the estimated power spectra are described by their covariance matrix ${\rm Cov}(\tilde{C}^{ab}_\ell,\tilde{C}^{cd}_{\ell'})$. As shown in \todo{cite}, a dominant component of the power spectrum covariance for many cosmological fields, including most large-scale structure tracers, is the so-called ``disconnected'' or ``Gaussian'' covariance, corresponding to the disconnected component of the fields' trispectrum (the only non-zero component for Gaussian random fields). The Gaussian pseudo-$C_\ell$ covariance is thus:
    \begin{equation}\nonumber
      {\rm Cov}(\tilde{C}^{ab}_\ell,\tilde{C}^{cd}_{\ell'})=\sum_{mm'}\frac{\langle \tilde{a}_{\ell m}\tilde{c}_{\ell'm'}^*\rangle\langle \tilde{b}^*_{\ell m}\tilde{d}_{\ell'm'}\rangle+(c\leftrightarrow d)}{(2\ell+1)(2\ell'+1)},
    \end{equation}
    where $(c\leftrightarrow d)$ denotes exchanging the roles of fields $c$ and $d$. 
    
    The key to estimating the Gaussian covariance is then the calculation of general two-point correlators of the form $\langle\tilde{a}_{\ell m}\tilde{c}_{\ell'm'}^*\rangle$. Using Eq. \ref{eq:mask_harmonic}, we can write this term as:
    \begin{equation}\label{eq:general_2pt}
      \langle\tilde{a}_{\ell m}\tilde{c}_{\ell'm'}^*\rangle=
      \sum_{\ell''m''}{\cal M}_{\ell m,\ell''m''}(w_a)\,{\cal M}^{*}_{\ell' m',\ell''m''}(w_c)C_{\ell''}^{ac}.
    \end{equation}
    The standard narrow kernel approximation (NKA) now proceeds by exploiting that, for reasonably broad masks, ${\cal M}_{\ell m,\ell''m''}(w)$ is narrowly peaked around $(\ell,m)=(\ell'',m'')$, and thus the power spectrum $C^{ac}_{\ell''}$ above may be simply approximated as $C^{ac}_{\ell''}\simeq(C^{ac}_\ell+C^{ac}_{\ell'})/2\equiv C^{ac}_{(\ell,\ell')}$. Combined with the following useful property of the map-level mode-coupling coefficients
    \begin{equation}
      \sum_{\ell''m''}{\cal M}_{\ell m,\ell''m''}(w_a)\,{\cal M}^{*}_{\ell' m',\ell''m''}(w_c)={\cal M}_{\ell m,\ell'm'}(w_aw_c),
    \end{equation}
    this then allows us to approximate the Gaussian covariance as:
    \begin{align}\nonumber
      {\rm Cov}(\tilde{C}^{ab}_\ell,\tilde{C}^{cd}_{\ell'})=
      &C^{ac}_{(\ell,\ell')}C^{bd}_{(\ell,\ell')}\Xi_{\ell\ell'}(w_aw_c,w_bw_d)\\\label{eq:nka_simple}
      &+(c\leftrightarrow d).
    \end{align}
    Calculating the Gaussian covariance thus reduces to estimating the same type of coupling coefficients required by the pseudo-$C_\ell$ estimator, this time involving the real-space product of pairs of masks in real space. Additionally, we require a reasonable guess for the true underlying power spectra of the fields involved (which, potentially, could be approximated from the measured spectra themselves).

    As described in \cite{2010.09717}, the accuracy of the NKA covariance can be further improved in two ways. The first important effect is the residual impact of mode coupling in the power spectrum before pulling it out of the sum in Eq. \ref{eq:general_2pt}. This is important for particularly complicated masks, where the core assumption of the NKA (the narrowness of the map-level kernels) is less accurate. This may be achieved by replacing the true power spectra in Eq. \ref{eq:nka_simple} (e.g. $C_\ell^{ac}$) by
    \begin{equation}\label{eq:clbar0}
      \overline{C}^{ab}_\ell\equiv\frac{\left\langle \tilde{C}^{ac}_\ell\right\rangle}{\langle w_aw_c\rangle_{\rm sky}},
    \end{equation}
    where the numerator is the mode-coupled spectrum (Eq. \ref{eq:pcl}), and the denominator is the sky average of the product of both masks
    \begin{equation}
      \langle w_aw_c\rangle_{\rm sky}\equiv\int\frac{d\nv}{4\pi}\,w_a(\nv)w_c(\nv).
    \end{equation}
    Here, $\langle\tilde{C}^{ab}_\ell\rangle$ may be calculated from a guess theory $C_\ell$ by convolving it with the pseudo-$C_\ell$ mode-coupling matrix, or it may be simply replaced by the pseudo-$C_\ell$ measured directly from the data.

    Secondly, all cosmological datasets are affected by some form of measurement noise (e.g. shape noise in cosmic shear), the properties of which (e.g. its variance) may vary across the sky. Approximating this noise to be uncorrelated between different pixels, we can calculate its impact on the general two-point correlators needed to calculate the Gaussian covariance (Eq. \ref{eq:general_2pt}) as
    \begin{align}\label{eq:general_2pt_wnoise}
      \langle\tilde{a}_{\ell m}\tilde{c}_{\ell'm'}^*\rangle
      =&\sum_{\ell''m''}{\cal M}_{\ell m,\ell''m''}(w_a)\,{\cal M}^{*}_{\ell' m',\ell''m''}(w_c)C_{\ell''}^{ac}\\\nonumber
      &+\delta_{ac}^K\,{\cal M}_{\ell m,\ell'm'}\left(\Omega_{\rm pix}\sigma_a^2w_a^2\right),
    \end{align}
    where the map $\sigma_a^2(\nv)$ denotes the noise variance in the pixel along direction $\nv$, and $\Omega_{\rm pix}$ is the pixel area. This additional contribution leads to the introduction of ``signal-signal'', ``signal-noise'', and ``noise-noise'' components of the total covariance matrix (see Eq. 2.29 of \cite{2010.09717}), all of which depend on pseudo-$C_\ell$ mode-coupling coefficients involving either the fields' masks or their mask-weighted noise variance maps.

    This improved NKA (iNKA) has been found to be remarkably accurate even for data with highly complex sky masks and inhomogeneous noise properties \todo{cites}.

  \subsection{Catalogue-based $C_\ell$s}\label{ssec:pcls.cat}
    Numerous cosmological datasets rely on the measurement of astrophysically-interesting quantities at the discrete positions of observed sources. Canonical examples of this are cosmic shear, based on the observed ellipticity of galaxies, fast-radio bursts, based on the observed dispersion measure (DM), or galaxy clustering, where the galaxy density is defined by the discrete positions of galaxies. In general such a catalogue-based field may be described in terms of its mask $w_a(\nv)$ and the masked map $\tilde{a}(\nv)$ as:
    \begin{align}\label{eq:mask_delta}
       &w_a(\nv)\equiv\sum_iw_i\,\delta^D(\nv,\nv_i),\\\label{eq:masked_cat_map}
       &\tilde{a}(\nv)=\sum_iw_ia_i\,\delta^D(\nv,\nv_i).
    \end{align}
    Here, all sums run over all sources in the catalogue, $w_i$ is the weight associated with the $i$-th source, $\nv_i$ are its sky coordinates, $a_i$ is the field value measured at that position, and $\delta^D(\nv_1,\nv_2)$ is the Dirac delta function defined on the sphere. The spherical harmonic transforms of the mask and masked map are then
    \begin{equation}
      w^a_{\ell m}=\sum_iw_i Y_{\ell m,i}^*,\hspace{12pt}
      \tilde{a}_{\ell m}=\sum_iw_ia_iY_{\ell m,i}^*,
    \end{equation}
    with $Y_{\ell m,i}\equiv Y_{\ell m}(\nv_i)$. The introduction of efficient algorithms for the calculation of discrete spherical harmonic transform on irregular grids \todo{cites}, makes it possible to calculate these spherical harmonic coefficients directly from the catalogues, circumventing the need to bin sources into regular finite pixel grids, and the numerical instabilities (e.g. aliasing, finite-resolution effects) that come with it. This may in fact lead to a computational speed-up with respect to the standard (map-based) pseudo-$C_\ell$ approach when targeting relatively large scales or sparse samples.

    As described in \cite{2407.21013}, the only additional step in the estimation of power spectra for catalogue-based fields, is to account for the impact of shot noise in the power spectra of both the mask and the masked field in the case of auto-correlation. Consider, for example, the mask power spectrum $\tilde{C}_\ell^w$: in the case of a continuous and reasonably simple mask, this power spectrum attains its largest values at low $\ell$, with the monopole $\ell=0$ being proportional to the total sky area covered, and then quickly drops and oscillates around zero. This decay is typically fast enough that one need only consider the mask power spectrum up to a modest maximum multipole $\ell_{\rm max}$, and the oscillatory behaviour may be tamed if needed by apodising the mask, for example. In the case of a catalogue-based mask, where sources are randomly sampled over similar sky area, the power spectrum is also largest at low $\ell$s, again tracking the total area covered by the sample, but then drops towards a constant shot-noise value $\tilde{N}_w$, given by the contribution from self-pairs to the power spectrum:
    \begin{equation}\label{eq:Nw}
      \tilde{N}_w\equiv\frac{1}{4\pi}\sum_iw_i^2.
    \end{equation}
    We show this explicitly in Fig. \todo{make figure}. Likewise, the power spectrum of the masked field reaches a noise-like floor $\tilde{N}_a$, given by the contribution from self-pairs:
    \begin{equation}\label{eq:Na}
      \tilde{N}_a\equiv\frac{1}{4\pi}\sum_iw_i^2a_i^2.
    \end{equation}
    As shown in \cite{2407.21013}, subtracting these noise-like components from the mask and masked field auto-spectra results in a numerically stable and unbiased estimator, which is also completely immune to any white noise component (homogeneous or otherwise).

    At this point, the challenges in adapting the iNKA approach to compute covariances of catalogue-based power spectra become apparent: the approximation involves calculating power spectra of map products (hidden inside the coupling coefficients in Eq. \ref{eq:nka_simple}), as well their sky average in Eq. \ref{eq:clbar0}. For catalogue-based masks (Eq. \ref{eq:mask_delta}), such products are either zero, in the case of cross-correlations between different catalogues, or they diverge for auto-correlations (since they involve products of coincident delta functions). Regularising this behaviour is the main topic of this paper.

\section{Analytical covariances for catalogue-based fields}\label{sec:theo}
  \subsection{Disjoint catalogues and cross-correlations}\label{ssec:theo.alldiff}
    For simplicity, let us start by considering a covariance matrix element of the form ${\rm Cov}(\tilde{C}^{ab}_\ell,\tilde{C}^{cd}_{\ell'})$, in which all fields involved, including their source catalogues, are different. Consider first the sky average $\langle w_aw_c\rangle_{\rm sky}$ required by the iNKA (Eq. \ref{eq:clbar0}). A priori, the product of both masks is zero, and thus so should $\langle w_aw_b\rangle_{\rm sky}$ be. Note, however, that we may write this quantity in terms of the pseudo-$C_\ell$ of both masks, using Parseval's theorem:
    \begin{equation}
      \langle w_aw_b\rangle_{\rm sky}=\sum_{\ell=0}^{\ell_{\rm max}}\frac{2\ell+1}{4\pi}\tilde{C}^{w_a,w_b}_\ell.
    \end{equation}
    As we argued in the last section, in general the $\tilde{C}^{w_aw_a}_\ell$ peaks at $\ell=0$, where it scales with the area over which both catalogues overlap, and then decays, eventually oscillating around zero (see Fig. \todo{ref}). The sum above should therefore be dominated by the lowest multipoles, and it only reaches zero in the limit $\ell_{\rm max}=\infty$ after all the incoherent oscillations in $\tilde{C}^{w_aw_b}_\ell$, amplified by the prefactor $2\ell+1$, cancel out. A similar expression appears when considering the numerator of Eq. \ref{eq:clbar0}:
    \begin{equation}
      \left\langle\tilde{C}^{ab}_\ell\right\rangle=\sum_{\ell''=0}^{\ell_{\rm max}}\frac{2\ell''+1}{4\pi}\tilde{C}^{w_aw_b}_\ell\,\Upsilon^{ab}_{\ell\ell'\ell''},
    \end{equation}
    where
    \begin{equation}
      \Upsilon^{ab}_{\ell\ell'\ell''}\equiv\sum_{\ell'}(2\ell'+1)C^{ab}_{\ell'}\wtj{\ell}{\ell'}{\ell''}{0}{0}{0}^2.
    \end{equation}
    Although both of these expressions are only strictly correct in the limit $\ell_{\rm max}\rightarrow\infty$, we can consider their behaviour as a function of $\ell_{\rm max}$. This is shown in Fig. \todo{fig}, together with their ratio $\overline{C}^{ab}_\ell$. As the figure shows, both quantities seemingly converge at $\ell_{\rm max}\ll\ell_{\rm ip}\equiv\pi/\theta_{\rm ip}$, where $\theta_{\rm ip}$ is the mean inter-source distance in the sample. Interestingly, their ratio is remarkably stable, even at very low values of $\ell_{\rm max}$. For values of $\ell_{\rm max}\gtrsim\ell_{\rm ip}$ the sums become more numerically unstable.

    Choosing a finite value for $\ell_{\rm max}$ (e.g. $\ell_{\rm ip}$) is equivalent to band-limiting the harmonic transform of the mask or, in other words, by treating each source not as a point, but as a ``cloud'' of finite size $\theta_{\rm c}$. For the purposes of covariance matrix estimation, we therefore begin by applying a Gaussian smoothing to the harmonic coefficients of the mask with standard deviation $\theta_{\rm c}$. I.e. we replace:
    \begin{align}\label{eq:wbar}
      w^a_{\ell m}\rightarrow \bar{w}^a_{\ell m}\equiv w^a_{\ell m}\phi^a_\ell,\hspace{6pt}
      \phi^a_\ell\equiv\exp\left[-\frac{\ell(\ell+1)\theta_{\rm c}^2}{2}\right].
    \end{align}
    In order to account for both the finite inter-particle distance of a given catalogue, as well as the possible band limit $\ell_{\rm max}$ imposed when analysing a given field, we use $\theta_c=\sqrt{\theta_{\rm ip}^2+(\pi/\ell_{\rm max})^2}$. Fig. \todo{fig} that this choice reproduces well the numerically stable value that the sums above converge to. Crucially, the replacement of delta function with Gaussian kernels addresses the other key problem of the iNKA approach when applied to catalogues. The resulting smoothed masks are defined continuously across the sky, allowing us to compute their real-space product in a numerically stable manner.

    From this discussion, we may now posit the following expression for the covariance of power spectra involving distinct catalogue-based fields:
    \begin{align}\nonumber
      {\rm Cov}(\tilde{C}^{ab}_\ell,\tilde{C}^{cd}_{\ell'})=
      &\overline{C}^{ac}_{(\ell,\ell')}\overline{C}^{bd}_{(\ell,\ell')}\Xi_{\ell\ell'}(\bar{w}_a\bar{w}_c,\bar{w}_b\bar{w}_d)\\\label{eq:nka_catdiff}
      &+(c\leftrightarrow d),
    \end{align}
    where $\bar{w}_a$ is the smooth mask given by Eq. \ref{eq:wbar}, and
    \begin{align}\nonumber
      &\overline{C}^{ab}_\ell\equiv\frac{\left\langle\tilde{C}^{ab}_\ell\right\rangle}{\langle \bar{w}_a\bar{w}_b\rangle_{\rm sky}},\\
      &\langle\bar{w}_a\bar{w}_b\rangle_{\rm sky}=\sum_{\ell}\frac{2\ell+1}{4\pi}\tilde{C}^{w_aw_b}_\ell\phi^a_\ell\phi^b_\ell. \label{eq:inka_cell}
    \end{align}

  \subsection{Auto-correlations}\label{ssec:theo.allsame}
    The case of covariance matrix elements involving the same catalogue require additional care. On the one hand, the shot noise contribution to the pseudo-$C_\ell$ of the masks and the masked map, caused by self-pairs, may lead to numerical instabilities. On the other hand, this shot noise does contribute to the statistical uncertainties, and must therefore be taken into account in the calculation of the covariance matrix.

    We start by considering the general two-point correlator $\langle\tilde{a}_{\ell m}\tilde{c}^*_{\ell'm'}\rangle$, first introduced in Eq. \ref{eq:general_2pt} as the building block of the Gaussian covariance, in the case of catalogue-based fields with $a=b$, separating the contribution from self-pairs:
    \begin{align}\nonumber
      \langle \tilde{a}_{\ell m}\tilde{a}^*_{\ell'm'}\rangle
      =&\sum_{i\neq j}w_iw_j\langle a_ia_j\rangle Y^*_{\ell m,i}Y_{\ell'm',j}\\\nonumber
      &+\sum_i w_i^2\langle a_i^2\rangle\,Y^*_{\ell m,i}Y_{\ell'm',i}\\\label{eq:general_2pt_cat}
      =&\sum_{i\neq j}(\cdots)+{\cal M}_{\ell m,\ell'm'}(\overline{(aw_a)^2}),
    \end{align}
    Where the last term is the map-level coupling coefficient (defined in Eq. \ref{eq:map_mcm}) for a map given by
    \begin{equation}\label{eq:self_var}
      \overline{(aw_a)^2}(\nv)\equiv\sum_i w^2_i\langle a_i^2\rangle\,\delta^D(\nv,\nv_i),
    \end{equation}
    which tracks the weighted variance of $a$ at the position of each source. Eq. \ref{eq:general_2pt_cat} is strongly reminiscent of the expression for the two-point correlator of a field with white noise, introduced in Eq. \ref{eq:general_2pt_wnoise}, with the self-pairs contribution playing the role of noise, and the sum over distinct pairs that of the signal.

    We can then estimate the contribution from different pairs (i.e. the first term in the previous equation) following the same rationale we used in the case of different-catalog covariances: we construct a mask by replacing each source with a Gaussian cloud, and use the square of this mask to construct the mode-coupling coefficients $\Xi_{\ell\ell'}$. In doing so, we must avoid the contribution to the square mask from self-pairs, however, since we are already accounting for this in the second term in Eq. \ref{eq:general_2pt_cat}. To calculate this contribution, let us write the harmonic coefficients of the square mask, separating out the contribution from self-pairs:
    \begin{align}
      (\bar{w}^2)_{\ell m}
      %&=\int d\nv\,Y^*_{\ell m}(\nv)\sum_i w_i\phi_a(\nv\cdot\nv_i)\,\sum_j w_j\phi_a(\nv\cdot\nv_j)\\
      %&=\sum_{i\neq j}(\cdots)+\sum_i w_i^2\,\int d\nv\,Y^*_{\ell m}(\nv)\left[\phi_a(\nv\cdot\nv_i)\right]^2\\
      &=\sum_{i\neq j}(\cdots)+\sum_i w_i^2\,(\phi^2_a)_\ell\,Y^*_{\ell m,i},
    \end{align}
    where $(\phi^2_a)_\ell$ is the harmonic transform of the square of $\phi_a$:
    \begin{align}
      (\phi^2_a)_\ell
      &\equiv 2\pi\int_{-1}^1 d\mu\,\phi_a(\mu)\,P_\ell(\mu)\\
      &=\frac{1}{4\pi\theta_{{\rm ip},a}^2}\exp\left[-\frac{1}{4}\ell(\ell+1)\theta_{{\rm ip},a}^2\right],
    \end{align}
    for a Gaussian kernel, where $P_\ell(\mu)$ are the Legendre polynomials. We can thus write the contribution to the squared mask from self-pairs as another catalog-based field, given by:
    \begin{equation}
      (\overline{w^2})_{\ell m}\equiv \sum_i w_i^2(\phi^2_a)_\ell\,Y^*_{\ell m,i}.
    \end{equation}

    With this, we can then write the iNKA version of the two-point correlator in Eq. \ref{eq:general_2pt} as
    \begin{equation}\label{eq:general_2pt_cat_inka}
      \langle \tilde{a}_{\ell m}\tilde{a}^*_{\ell'm'}\rangle\simeq \overline{C}^{aa}_{(\ell,\ell')}{\cal M}_{\ell m,\ell'm'}(\bar{w}^2-\overline{w^2})+{\cal M}_{\ell m,\ell'm'}(\overline{(aw)^2}).
    \end{equation}
    Here, we must remember to subtract all self-pair contributions to $\overline{C}^{aa}_\ell$, since these are also contained in the second term. I.e.
    \begin{align}
      &\overline{C}^{aa}_\ell\equiv\frac{\tilde{C}^{aa}_\ell-\tilde{N}_a}{\langle \bar{w}^2-\overline{w^2}\rangle},\\
      &\langle \bar{w}^2-\overline{w^2}\rangle=\sum_\ell\frac{2\ell+1}{4\pi}(\tilde{C}^{ww}_\ell-\tilde{N}_w)(\phi^a_\ell)^2,
    \end{align}
    with $\tilde{N}_a$ and $\tilde{N}_w$ given in Eqs. \ref{eq:Na} and \ref{eq:Nw}, respectively. As we said before, the structure of Eq. \ref{eq:general_2pt_cat_inka} is similar to that of a field with inhomogeneous white noise, as described in \cite{2010.09717}, with the self-pair contribution playing the role of the inhomogeneous noise component. We can thus calculate the power spectrum covariance for catalogue-based fields by including signal-noise and noise-noise terms that account for the shot noise contribution from self-pairs.

    It is important to note that the parallel between the structure described here and the signal/noise contributions described in \cite{2010.09717} is more than a simple coincidence: the noise-like contribution from self-pairs is an intrinsic aspect of fields sampled at discrete points that penalises sparser samples with larger statistical uncertainties.

  \subsection{General covariance structure}\label{ssec:theo.general}
    Putting together the results presented in the preceding two sections, we can write a general expression for the covariance matrix of pseudo-$C_\ell$s of catalog-based fields:
    \begin{widetext}
    \begin{align}\nonumber
      {\rm Cov}(\tilde{C}^{ab}_\ell,\tilde{C}^{cd}_{\ell'})
      =&\,\overline{C}^{ac}_{(\ell,\ell')}\overline{C}^{bd}_{(\ell,\ell')}\Xi_{\ell\ell'}\left(\bar{w}_a\bar{w}_c-\delta^K_{ac}\overline{w_a^2},\bar{w}_b\bar{w}_d-\delta^K_{bd}\overline{w_b^2}\right)+\delta^K_{bd}\,\overline{C}^{ac}_{(\ell,\ell')}\Xi_{\ell\ell'}\left(\bar{w}_a\bar{w}_c-\delta^K_{ac}\overline{w_a^2},\overline{(bw_b)^2}\right)\\\label{eq:main_result}
      &+\delta^K_{ac}\,\overline{C}^{bd}_{(\ell,\ell')}\Xi_{\ell\ell'}\left(\overline{(aw_a)^2},\bar{w}_b\bar{w}_d-\delta^K_{bd}\overline{w_b^2}\right)+\delta^K_{ac}\,\delta^K_{bd}\Xi_{\ell\ell'}\left(\overline{(aw_a)^2},\overline{(bw_b)^2}\right)\\\nonumber
      &+(c\leftrightarrow d),
    \end{align}
    \end{widetext}
    where the different ingredients entering this expression are:
    \begin{align}
      &\overline{C}^{ab}_\ell\equiv\frac{\tilde{C}^{ab}_\ell-\delta_{ab}\tilde{N}_a}{\sum_{\ell'}\frac{2\ell'+1}{4\pi}(\tilde{C}^{w_aw_b}_{\ell'}-\delta_{ab}\tilde{N}_w)\phi^a_\ell\phi^b_\ell},\\
      &(\bar{w}_a)_{\ell m}=\sum_iw^a_i\phi^a_\ell\,Y^*_{\ell m,i},\\
      &(\overline{w_a^2})_{\ell m}\equiv\sum_i(w^a_i)^2(\phi^2_a)_\ell Y^*_{\ell m,i},\\\label{eq:awa2}
      &\left(\overline{(aw_a)^2}\right)_{\ell m}\equiv\sum_i(a_iw^a_i)^2Y_{\ell m,i}^*,\\
      &\phi^a_\ell\equiv e^{-\frac{1}{2}\ell(\ell+1)\theta^2_{{\rm ip},a}},
      \hspace{12pt}(\phi^2_a)_\ell\equiv\frac{e^{-\frac{1}{4}\ell(\ell+1)\theta^2_{{\rm ip},a}}}{4\pi\theta^2_{{\rm ip},a}}.
    \end{align}
    $\tilde{N}_w$ and $\tilde{N}_a$ are given in Eqs. \ref{eq:Nw} and \ref{eq:Na}, respectively, and the coupling coefficients $\Xi_{\ell\ell'}(m_1,m_2)$ are defined in Eq. \ref{eq:Xi}.

    A final note must be made concerning the calculation of the noise-noise component of the Gaussian covariance (the last term in Eq. \ref{eq:main_result}). This term involves the construction of the local variance map in Eq. \ref{eq:self_var}, which requires an estimate of the field variance at each source. This is typically not immediately available, and must be estimated directly from the data. A na\"ive ansatz, which we use in Eq. \ref{eq:awa2}, is to simply replace $\langle a_i^2\rangle\rightarrow a_i^2$ for each source. However, the noise-noise contribution requires us to calculate the pseudo-$C_\ell$ of this map, which then involves the fourth-order moment of $a_i$ under this approximation. In other words, the fact that $\langle a_i^2\rangle\langle a_j^2\rangle\neq\langle a_i^2a_j^2\rangle$ leads to an over-estimation of the power spectrum of the local variance map at high-ell, which we must correct for. Under the assumption of Gaussian statistics this can be done analytically, by making use of Wick's theorem (for more details, see Appendix \ref{app:awa2_shotnoise}). 
    
    Furthermore, to simplify the discussion, the derivation above has assumed that we are interested in the covariance matrix of the pseudo-spectrum $\tilde{C}^{ab}_\ell$. In reality, as we discussed in Section \ref{ssec:pcls.cat}, we must always subtract the contribution from self-pairs $\tilde{N}_a$ (Eq. \ref{eq:Na}) in catalog-based fields, to ensure the estimator is unbiased and numerically stable. Thus, more precisely, we are in fact interested in the covariance of the combination $\tilde{C}^{ab}_\ell-\delta_{ab}^K\tilde{N}_a$, which leads to a similar correction to the pseudo-$C_\ell$ of the local variance map. In summary, any occurrence of the auto-spectrum of $\overline{(aw_a)^2}$ in the calculation must have the following additive constant offset subtracted:
    \begin{equation}\label{eq:4point_corr_spin0}
      \tilde{N}_{\overline{(aw_a)^2}}\equiv\frac{1}{4\pi}\sum_iw_i^4a_i^4.
    \end{equation}
    Note that, just like $\tilde{N}_a$ is the contribution to the the auto-correlation $\tilde{C}^a_\ell$ from self-pairs, so is $\tilde{N}_{\overline{(aw_a)^2}}$ the contribution to the pseudo-spectrum of $\overline{(aw_a)^2}$ from self-pairs. A proof of this can be found in Appendix \ref{app:awa2_shotnoise}.

  \subsection{Momentum fields and galaxy clustering}\label{ssec:theo.clust}
    There is a subclass of catalogue-based fields, called ``momentum'' fields in this paper, which represent cosmological tracers weighted by the local density of galaxies. An example is the radial momentum field $\pi_r\sim (1+\delta_g)\,v_r$, used in the measurement of the low-redshift kinematic Sunyaev-Zel'dovich effect \cite{2512.14625,2604.19744}, in which the quantity weighted by the local number density is the spin-0 radial velocity field. Momentum fields may also be constructed from quantities with arbitrary spin, and can be exploited in e.g. the study of the moving lens effect, or in general directional stacking analyses \cite{2605.15947,2605.18938}. In these cases, the masked field takes the same form as that of the standard sampled field, Eq. \ref{eq:masked_cat_map}, where $a_i$ is the value of a spin-$s$ field at the position of the $i$-th source. However, unlike the case of sampled fields, the mask of a momentum field represents the \emph{expected} mean number density of sources (i.e. the number that would have been observed in an otherwise homogeneous universe), as opposed to the actual observed density. This is often called the ``completeness'' or ``depth'' map, and may be represented either by a continuous map, or by a set of random points tracking the local expected density of sources. In the first case (that of a continuous mask), Eq.~\ref{eq:main_result} is modified by replacing the signal-signal covariance kernel by its standard iNKA expression, as shot noise is not an issue. 

    In turn, a randoms-based mask takes the form
    \begin{equation}
      w(\nv)=\sum_{r\in R} \alpha\,w_r\,\delta^D(\nv,\nv_r),
    \end{equation}
    where we have specified that the sums runs only over the random points. The normalisation constant $\alpha\equiv\sum_{d=D}w_d/\sum_{r=R}w_r$ ensures that $w$ represents the local density of sources in the data. For simplicity of notation in what follows we will absorb $\alpha$ into the random weights $w_i$ (i.e. they are normalised such that $\sum_{r\in R}w_r=\sum_{d\in D}w_d$).
    
    Galaxy clustering can itself be treated as momentum field with the following minimal modifications:
    \begin{enumerate}
      \item Each source in the data is assigned a field value of 1.
      \item The mean density (estimated from the mask) is subtracted from the masked field.
    \end{enumerate}
    I.e., in this case the masked field is
    \begin{align}\nonumber
      \delta_g(\nv)
      &=\sum_{d\in D}w_d\,\delta^D(\nv,\nv_d)-w(\nv)\\\nonumber
      &=\sum_{d\in D}w_d\,\delta^D(\nv,\nv_d)-\sum_{r\in R}w_r\,\delta^D(\nv,\nv_r),
    \end{align}
    where, in the second line, we have assumed that the mask is defined in terms of randoms. Since the random catalogue also contributes to the shot noise of the masked map (i.e. the contribution from self-pairs), the noise terms in the covariance matrix must also include this contribution. I.e. the map used to calculate the noise-like coupling coefficients in the case of galaxy clustering should be
    \begin{equation}
      \overline{(aw_a)^2}(\nv)=\sum_{d\in D}w_d^2\,\delta^D(\nv,\nv_d)+\sum_{r\in R}w_r^2\,\delta^D(\nv,\nv_r),
    \end{equation}
    instead of Eq. \ref{eq:self_var}. Likewise, the shot-noise contribution to the pseudo-spectrum must account for randoms:
    \begin{equation}
      \tilde{N}_a=\frac{1}{4\pi}\left(\sum_{d\in D}w_d^2+\sum_{r\in R}w_r^2\right)
    \end{equation}
    for galaxy clustering. Finally, no additive correction to the noise-noise covariance $\Delta\tilde{C}^{\overline{(aw_a)^2}}_\ell$ (Eq.~\ref{eq:4point_corr_spin0}) is needed, as there is no fourth-order term $\langle a_i^2 a_j^2 \rangle$ in the case of clustering. These modifications do not apply to other types of momentum fields, in which the mean density contribution from randoms is not subtracted from the masked map. However, in any momentum field using randoms, all mask elements entering the Gaussian covariance in Eq. \ref{eq:main_result} (e.g. $\bar{w}_a$, $\overline{w_a^2}$, $\tilde{N}_w$, $\tilde{C}^{w_aw_b}_\ell$), as well as the four-point additive correction (Eq.~\ref{eq:4point_corr_spin0}) must be calculated from the randoms, and not the data.

  \subsection{Exact summation}\label{ssec:theo.brute}
For sufficiently sparse samples, it is feasible to evaluate the covariance of the decoupled power spectra exactly, without resorting to the iNKA. We describe this calculation here for the auto-correlation of a single catalogue-based scalar field $a$. The decoupled estimator $\hat{C}_\ell^{aa}$ is obtained from the pseudo-$C_\ell$ (Eq.~\ref{eq:pcl_def}) by subtracting the shot-noise bias $\tilde{N}_a$ (Eq.~\ref{eq:Na}) and applying the inverse of the mode-coupling matrix (Eq.~\ref{eq:pcl}),
\begin{equation}\label{eq:invert}
    \hat{C}_{\ell}^{aa} = \sum_{\ell_1} [M^{-1}]_{\ell \ell_1} \, \big(\tilde{C}_{\ell_1}^{aa} - \tilde{N}_a\big).
\end{equation}
Inserting this into the covariance and expanding the product gives
\begin{equation} \label{eq:cov2}
\begin{split}
\mathrm{Cov}&(\hat{C}_\ell^{aa}, \hat{C}_{\ell'}^{aa}) =  \sum_{\ell_1,\ell_2} [M^{-1}]_{\ell\ell_1} [M^{-1}]_{\ell_2\ell'} \\ &\times  \Big[ \left\langle \tilde{C}_{\ell_1}^{aa} \tilde{C}_{\ell_2}^{aa} \right\rangle - \left\langle \tilde{N}_a\,\tilde{C}_{\ell_1}^{aa} \right\rangle 
- \left\langle \tilde{N}_a\,\tilde{C}_{\ell_2}^{aa} \right\rangle + \left\langle \tilde{N}_a^2\right\rangle\Big] \\ & - \left\langle \hat{C}_\ell^{aa} \right\rangle \left\langle \hat{C}_{\ell'}^{aa} \right\rangle,
\end{split}
\end{equation}
where $\ell_1$ and $\ell_2$ are summed multipoles.
 
Under the assumption of Gaussianity, each correlator in Eq.~\ref{eq:cov2} reduces by Wick's theorem to products of the field two-point function evaluated at pairs of sources. Splitting the field value at the $i$-th source into signal and measurement-noise contributions, $a_i = s_i + n_i$, this two-point function reads
\begin{equation} \label{eq:field_corr}
    \left\langle a_i\,a_j\right\rangle = \left\langle s_i\,s_j\right\rangle + \sigma_{N}^2\, \delta^K_{ij},
\end{equation}
with uncorrelated noise $\langle n_i n_j\rangle = \sigma_{N}^2 \delta^K_{ij}$, and per-source signal variance
\begin{equation}
    \sigma_{S}^{2} \equiv \langle s_{i}^2\rangle = \sum_\ell \frac{2\ell + 1}{4\pi}\, \hat{C}_\ell^{aa}.
\end{equation}
The noise variance $\sigma_{N}^2$ is fixed by contracting both sides of Eq.~\ref{eq:field_corr}.
 
Carrying out the sums over $m$ with the spherical-harmonic addition theorem, $\sum_m Y_{\ell m, i}^* Y_{\ell m, j} = (2\ell + 1)\,P_{\ell, ij}/4\pi$, where $P_{\ell, ij} \equiv P_\ell(\nv_i \cdot \nv_j)$ and $P_\ell$ are the Legendre polynomials, and using the fact that the estimator is unbiased, so that $\langle \hat{C}_\ell^{aa}\rangle\langle \hat{C}_{\ell'}^{aa}\rangle = C_\ell^{aa}\,C_{\ell'}^{aa}$, the covariance becomes
\begin{widetext}
\begin{equation} \label{eq:cov_brute_long}
\begin{split}
    \mathrm{Cov}(\hat{C}_\ell^{aa}, \hat{C}_{\ell'}^{aa}) = & \sum_{\ell_1,\ell_2} [M^{-1}]_{\ell\ell_1} [M^{-1}]_{\ell_2\ell'} \Bigg[\langle \tilde{C}_{\ell_1}^{aa} \rangle \langle\tilde{C}_{\ell_2}^{aa} \rangle
 \\ + & \sum_{i,j, p, q}  \frac{w_i w_jw_p w_q}{(4\pi)^2} P_{\ell_1, ij}   P_{\ell_2, pq} \left\langle a_i a_p \right\rangle \left\langle a_j a_q \right\rangle 
+ \sum_{i,j, p, q}  \frac{w_i w_j w_p w_q}{(4\pi)^2} P_{\ell_1, ij} P_{\ell_2, pq}
  \left\langle a_i a_q \right\rangle \left\langle a_j a_p \right\rangle\\
-&\; \langle \tilde{N}_a \rangle \langle\tilde{C}_{\ell_1}^{aa} \rangle-  \sum_{i, j, k} \frac{2 w_k^2 w_i\,w_j}{(4\pi)^2} P_{\ell_1, ij}
 \left\langle a_i\,a_k\right\rangle \left\langle a_j\,a_k\right\rangle 
- \langle \tilde{N}_a \rangle \langle\tilde{C}_{\ell_2}^{aa} \rangle-  \sum_{i, j, k} \frac{2 w_k^2 w_i\,w_j}{(4\pi)^2} P_{\ell_2, ij}
 \left\langle a_i\,a_k\right\rangle \left\langle a_j\,a_k\right\rangle \\
+& \left(\tilde{N}_{w} \, \sigma_{S}^{2}\right)^2 +  \sum_{ij}  \frac{2 w_i^2 w_j^2 }{(4\pi)^2} \left\langle a_i\,a_j\right\rangle^{2} \Bigg]
- C_\ell^{aa}\,  C_{\ell'}^{aa}.
\end{split}
\end{equation}
\end{widetext}
 
This expression simplifies considerably. The two signal--signal terms (the second and third) are identical under relabelling of the dummy indices. The expectation value of the pseudo-$C_\ell$ is the coupled signal spectrum plus the self-pair shot-noise floor, which splits into a field-variance and a measurement-noise contribution,
\begin{equation}
    \left\langle \tilde{C}_{\ell}^{aa} \right\rangle = \sum_{\ell_1} M_{\ell\ell_1} C_{\ell_1}^{aa}  + \tilde{N}_{w} \, \sigma_{S}^{2} + \tilde{N}_{w} \, \sigma_{N}^{2},
\end{equation}
with $\tilde{N}_w$ the mask shot noise (Eq.~\ref{eq:Nw}). Finally, separating each source sum into its equal- and distinct-index parts,
\begin{equation}
\begin{aligned}
    \sum_{i,j, p, q} &= \sum_{i=j, p= q} + \sum_{i=j, p\not= q} + \sum_{i\not=j, p= q} + \sum_{i\not=j, p\not= q},\\
    \sum_{i,j, k} &= \sum_{i=j, k} + \sum_{i\not=j, k},
\end{aligned}
\end{equation}
the equal-index (self-pair) contributions are exactly those removed by the debiasing in Eq.~\ref{eq:invert} (see also Appendix \ref{app:awa2_shotnoise}), while the remaining disconnected pieces cancel against $C_\ell^{aa}C_{\ell'}^{aa}$. The covariance therefore reduces to
\begin{equation} \label{eq:cov_brute}
\begin{split}
    \mathrm{Cov}&(\hat{C}_\ell^{aa}, \hat{C}_{\ell'}^{aa})  =  2\sum_{\ell_1,\ell_2} [M^{-1}]_{\ell\ell_1} [M^{-1}]_{\ell_2\ell'} \\
    &\times \sum_{i \not= j,\, p \not= q}\frac{w_i w_jw_p w_q}{(4\pi)^2}   P_{\ell_1, ij} P_{\ell_2, pq}
     \left\langle a_i a_p \right\rangle \left\langle a_j a_q \right\rangle.
\end{split}
\end{equation}
The equivalence of this expression with Eq.~\ref{eq:main_result} is shown in Appendix \ref{app:brute_nka}.

 Eq.~\ref{eq:cov_brute} is computationally prohibitive for large catalogues, as it scales with the number of sources $N_\mathrm{sources}$ and the maximum value of the multipoles $\ell_\mathrm{max}$  as $\mathcal{O}(N^4 \ell_\mathrm{max}^4)$. In addition, memory restrictions may apply depending on the available hardware. We therefore implement it using the \texttt{jax} \cite{jax2018github} library, with details given in Appendix \ref{app:brute_code}.
 

\section{Validation}\label{sec:val}
  We now turn to validating the analytical covariance formalism developed in Sect.~\ref{sec:theo} against sample covariance estimators using simulated data. We start with a suite of noiseless, low-resolution sky simulations covering all types of fields occurring in a multi-tracer analysis, before moving on to more realistic mock data tests mimicking sparse spin-0 catalogs (fast radio bursts), dense spin-2 catalogs (galaxy shear), and galaxy clustering.
  \subsection{General method validation}\label{ssec:val.general}
    We now validate our method against simulation-based covariances using a general noiseless simulation setup. We simulate Gaussian fields on the sky, masking it using the {\it Quaia} survey mask $w_s$ \todo{cite} using a \texttt{HEALPix} pixelization at resolution $N_{\rm side}=256$. We then simulate four mock catalogs as described below, and compute the analytical covariance using the simulation ground-truth $C_\ell$s as the input for the iNKA $\bar{C}_\ell^{ab}$ (Eq.~\ref{eq:inka_cell}) to avoid realization-dependent bias. Then we draw 1000 sky realizations, generate the set of 4 catalog fields, estimate their catalog-based power spectra using the method described in \todo{cites}, and compute their sample covariance. For the four fields that form the legs of the covariance, we consider galaxy clustering fields, momentum catalog fields, and tracer catalog fields, the latter two of which can also vary in spin. Furthermore, we may choose for some, or all, legs to be identical (same mock survey) or differ in catalog positions. We show the complete set of test configurations in Table \ref{tab:validation}. The three field types are generated in slightly different ways, as described hereinafter.
    \begin{itemize}
        \item {\it Tracer catalog:} We draw the field $a_{\ell m}$ as a Gaussian realization of a power-law power spectrum $C_\ell\propto\ell^\alpha$ with index $\alpha=-0.7$. We tested different values of $\alpha$ and verified that the results presented here are insensitive to this choice. In the case of spin 2, we draw uncorrelated $E$ and $B$ mode realizations. We draw one million source positions as inhomogeneous Poisson realizations with modulating mask $w_s$, and evaluate the field values at the exact positions using an irregular-grid inverse SHT \todo{cite}. We then construct the field calling the \texttt{nmt.NmtFieldCatalog} class with the positions and field values, and using $\ell_{\rm max}=3N_{\rm side}-1$. We use unit weights for simplicity.
        
        \item {\it Galaxy clustering:} We draw a Gaussian overdensity map realization $\delta_g$ at $N_{\rm side}=256$, using an input power spectrum $C_\ell^\delta\equiv A_n\, C^0_\ell$, with $C_\ell^0=1/(10+\ell)^{1.2}$ and $A_n^{-1}=n^2\sum_{\ell'}(2\ell'+1)C_{\ell'}^0/(4\pi)$, ensuring that the overdensity field has a standard deviation of $1/n$. We conservatively choose $n=5$ to ensure that the field is positive definite (i.e. $\delta_g>-1$) everywhere. We then draw mock sources as realizations of an inhomogeneous Poisson distribution with a modulating mean density proportional to $w_s(1+\delta_g)$. We also generate random catalogues using the modulating mask $w_s$. As written in Table~\ref{tab:validation}, the number of sources varies between \todo{update} $10^6$ and $2\times10^7$, in order to reduce the numerical noise from the simulation-based reference covariances. We then use \texttt{nmt.NmtFieldCatalogClustering} to construct a field with the corresponding source and random positions, or alternatively, we use $w_s$ as a completeness mask.
        
        \item {\it Momentum catalog:} We repeat the same steps as for galaxy clustering until generating the randoms, using one or ten million artificial sources and 20 million randoms. Similar to galaxy clustering, the number of sources is chosen such that the empirical covariance is sufficiently numerically stable to be compared with the analytical covariance. We then draw the spherical harmonic coefficients $v_{\ell m}$ of a ``velocity'' field from a power spectrum $C_\ell^v=1/(10+\ell)^{2.2}$, and evaluate the field's value at the exact catalog positions. We construct the momentum field using the \texttt{nmt.NmtFieldCatalogMomentum} class applied to the source positions, the corresponding velocity field values, and the random positions. Finally, we obtain the exact expression for the expected power spectrum of $\pi=v\,(1+\delta)$ from $C_\ell^\delta$ and $C_\ell^v$ (see Appendix A of \cite{2512.14625}).
    \end{itemize}

    \begin{table*}[t]
    \caption{\textbf{Validation cases for the iNKA covariance estimator.} The different columns are: the field types that form the 4 legs of the covariance, their spin modes, the number of mock sources in the catalog fields, the number of randoms (if randoms are being used for any of the 4 fields), the input power spectrum used to generate the sky modes, and the Kolmogorov-Smirnov PTEs assessing statistical correspondence  between empirical and analytical covariances for same-sources and different-sources (in brackets) scenarios.}
    \centering
    \scriptsize
    \setlength{\tabcolsep}{3pt}
    \renewcommand{\arraystretch}{1.12}
    
    \begin{tabular}{@{} l l l l l p{0.06\textwidth} p{0.06\textwidth}@{}}
    \toprule
    Case
    & Spin modes
    & $N_{\rm sources}$
    & $N_{\rm randoms}$
    & Power spectrum
    & KS PTE (same sources)
    & KS PTE (diff. sources)\\
    \midrule
    
    (catalog, spin-0)$^4$
    & $0000$
    & $2\times 10^7$
    & -
    & $1/(10+\ell)^{0.7}$
    & 0.45 & 0.43 \\

    (catalog, spin-2)$^4$
    & $EEEE$
    & $10^6$
    & -
    & $1/(10+\ell)^{0.7}$
    & 0.63 & 0.93 \\

    (catalog, spin-2)$^4$
    & $EBEB$
    & $10^6$
    & -
    & $1/(10+\ell)^{0.7}$
    & 0.06 & 0.28 \\

    (catalog, spin-2)$^4$
    & $BBBB$
    & $10^6$
    & -
    & $1/(10+\ell)^{0.7}$
    & 0.56 & 0.50 \\

    (catalog)$^4$, $00\times22$ 
    & $00EE$
    & $10^6$
    & -
    & $1/(10+\ell)^{0.7}$
    & 0.20 & 0.83 \\
    
    (catalog)$^4$, ($02$)$^2$
    & $0E\,0E$
    & $10^6$
    & -
    & $1/(10+\ell)^{0.7}$
    & 0.51 & 0.71 \\
    
    (catalog)$^4$, $22\times02$
    & $EE\,0E$
    & $10^6$
    & -
    & $1/(10+\ell)^{0.7}$
    & 0.79 & 0.63 \\
    
    (catalog)$^4$, $00\times02$
    & $000E$
    & $10^6$
    & -
    & $1/(10+\ell)^{0.7}$
    & 0.81 & 0.54  \\
    
    (catalog $\times$ map)$^2$
    & $0000$
    & $10^6$
    & -
    & $1/(10+\ell)^{0.7}$
    & 0.71 & 0.97 \\
    
    (clustering $\times$ catalog)$^2$
    & $0000$
    & $10^6$
    & $2\times 10^7$
    & $1.1\times10^{-3}/(10+\ell)^{1.2}$
    & 0.75 & \textcolor{red}{$0.00015$}$^\dagger$\\
    
    (clustering $\times$ map)$^2$
    & $0000$
    & $10^6$
    & -
    & $1.1\times10^{-3}/(10+\ell)^{1.2}$
    & 0.94 & 0.42 \\
    
    (clustering)$^4$
    & $0000$
    & $2\times 10^7$
    & $2\times 10^7$
    & $\delta_g$: $1.1\times10^{-3}/(10+\ell)^{1.2}$  % $3.7\times10^{-4}/(10+\ell)$
    & 0.12 & \textcolor{red}{0.003}$^\dagger$\\
    
    (momentum $\times$ map)$^2$
    & $0000$
    & $10^7$
    & $2\times 10^7$
    & $\delta_g$: $1.1\times10^{-3}/(10+\ell)^{1.2}$ %; $v_r$: $1/(10+\ell)^{2.2}$
    & 0.54 & 0.05$^\dagger$ \\

    (momentum, spin-2)$^4$ ($EE$)$^2$
    & $EEEE$
    & $10^6$
    & $2\times 10^7$
    & $\delta_g$: $1.1\times10^{-3}/(10+\ell)^{1.2}$
    & 0.11$^\dagger$ &  0.13\\

    (momentum, spin-2)$^4$
    & $EBEB$
    & $10^6$
    & $2\times 10^7$
    & $\delta_g$: $1.1\times10^{-3}/(10+\ell)^{1.2}$
    & 0.95 & \textcolor{red}{0.98}$^\dagger$\\

    (momentum, spin-2)$^4$
    & $BBBB$
    & $10^6$
    & $2\times 10^7$
    & $\delta_g$: $1.1\times10^{-3}/(10+\ell)^{1.2}$
    & \textcolor{red}{0.04} & 0.37$^\dagger$\\

    \bottomrule
    \end{tabular} \label{tab:validation} \\
    \footnotesize{$^\dagger$ The $\chi^2$ PTEs were computed with $\ell\in[2, 500]$ due to incomplete convergence of the empirical covariance.}
    \end{table*}

    As indicated in Table~\ref{tab:validation}, we separate the validation suite into three categories: scalar fields, spin-2 fields, and combinations thereof.  Figures \ref{fig:var_spin0}--\ref{fig:var_spin2xspin0} show how the simulated and analytical variance (the diagonal of the covariance matrix) compare across validation setups.
    
     For the covariance validation, we compute chi-squared statistics, $\chi^2 = X^T \Sigma Y^T$, where $X$ and $Y$ are the binned simulated power spectrum, $\langle\cdot\rangle$ is the mean over 1000 simulations, and $\Sigma=\text{Cov}(X,\, Y)$ is either the analytic or the sample covariance from the same set of simulations. The sample $\chi^2$ distribution is considered the benchmark.Figures \ref{fig:chi2_spin0_alldiff}--\ref{fig:chi2_spin2xspin0_allsame} show the $\chi^2$ distribution of all 1000 simulations for the full validation suite, distinguishing between the case when the catalogs entering the covariance have all different sources, and when they are all the same. We also plot a theoretical $\chi^2$ distribution with $n_{\rm bin}=20$ degrees of freedom, which the sample distribution is expected to converge to if $X$ is identical to $Y$. Where this is not the case, the (non analytical) theoretical distribution is omitted. We pass all $\chi^2$ tests, meaning that the distribution of probabilities for $\chi^2_{\rm sampled}$ to exceed $\chi^2_{\rm iNKA}$ ($\chi^2$ PTE) are compatible with a uniform distribution in all cases, measured by a Kolmogorov-Smirnov (KS) test. We define ``pass'' if the KS PTE is between 0.05 and 0.95. The KS PTEs are shown in Table~\ref{tab:validation}. Figure \ref{fig:chi2_overview} shows, for six representative cases, the $\chi^2$ PTE distribution in the small inset plots. %The galaxy clustering validation is barely passing, since discretely sampled overdensity drawn from a red power spectrum require very large sample sizes to avoid sampling noise at the smallest scales. In our concrete case, we increased the source count to 20 million to pass the $\chi^2$ test.
    \todo{Repeat the tests and comment on the outcome.}

    \begin{figure}
      \centering
      \includegraphics[width=0.49\textwidth]{figures/var_spin0_paper.pdf}
      \caption{Power spectrum variance estimates for spin-0 field validation cases, comparing the iNKA approximation (lines) against the sample variance from 1000 simulations (markers). Each covariance shown either has catalogs share the same sources (solid lines, dots), or have no sources in common (dashed lines, crosses). The cases are offset by factors of 10 for visual purposes.}\label{fig:var_spin0}
    \end{figure}

    \begin{figure}
      \centering
      \includegraphics[width=0.49\textwidth]{figures/var_spin2_paper.pdf}
      \caption{Same as Fig.~\ref{fig:var_spin0}, but showing the spin-2 validation cases.}\label{fig:var_spin2}
    \end{figure}

    \begin{figure}
      \centering
      \includegraphics[width=0.49\textwidth]{figures/var_spin2xspin0_paper.pdf}
      \caption{Same as Fig.~\ref{fig:var_spin0}, but showing the spin-0 $\times$ spin-2 validation cases.}\label{fig:var_spin2xspin0}
    \end{figure}
    
    \begin{figure*}
      \centering
      \includegraphics[width=\textwidth]{figures/chi2_overview_paper.pdf}
      \caption{$\chi^2$ distributions for the iNKA covariance and the empirical covariance for a selection of validation cases, with all catalog fields sharing all the same sources. Small inset plots show the distribution of $\chi^2_{\rm iNKA}$ PTEs with respect to $\chi^2_{\rm sim}$.}\label{fig:chi2_overview}
    \end{figure*}



  \subsection{Sparse catalogues: FRBs}\label{ssec:val.sparse}
    A typical example for a very sparse tracer are Fast Radio Bursts (FRBs). Currently a couple of thousand FRB sources are known with by far the largest number having been detected by the Canadian Hydrogen Intensity Mapping Experiment
 \citep[CHIME,][]{2026ApJS..283...34C}. While we expect this number to grow significantly over the next decade, the corresponding catalogues will overall stay sparse with expected number densities $\approx 10^{-3}\;\mathrm{arcmin}^{-2}$.

  \subsection{Dense catalogues: cosmic shear}\label{ssec:val.dense}
    \begin{figure*}
      \centering
      \includegraphics[width=0.49\textwidth]{figures/cov_shear_EE.pdf}
      \includegraphics[width=0.49\textwidth]{figures/cov_shear_BB.pdf}
      \caption{Power spectrum covariance for the DES-like simulated cosmic shear sample. Results are shown for the covariance of the $E$-mode and $B$-mode power spectra (left and right panels). The blue lines and points show results for simulations in which only the lensing shear signal is sampled at the positions sources, whereas results including both signal and Gaussian uncorrelated noise are shown in red. In the latter case we assume a noise standard deviation per component of $\sigma_e=0.27$. The sample variance from 1000 realisation is shown as points, with the analytical prediction as solid lines. The dashed line and crosses show the first off-diagonal element of the covariance matrix.}
      \label{fig:cov_shear}
    \end{figure*}
    In this second example we consider the opposite regime to that studied in the previous section: a high-density sample ($\bar{n}\sim O(1-10\,{\rm arcmin}^{-2})$) where power spectra can be recovered up to very high multipoles, and where the field values sampled at any individual source are noise-dominated. These are the characteristics of cosmic shear data, a key cosmological observable, and one where catalogue-based power spectrum methods are particularly valuable \cite{2407.21013,2605.13543}. In this regime, a brute-force approach to the covariance matrix is computationally intractable, and the analytical approximations described here become most useful.

    In order to test our approach in a realistic setting, we use the positions of galaxies from the Dark Energy Survey (DES) Year-3 (Y3) cosmic shear sample \todo{cite}. Specifically, we select galaxies from the third redshift bin used in the analysis of \todo{ctie}. We generate Gaussian realisations of the cosmic shear signal and sample it at the positions of these sources. We do so by generating a Gaussian realisation of the $E$-mode harmonic coefficients of the cosmic shear field, sampled from the theoretical shear power spectrum for this sample, calculated assuming a {\sl Planck}-like cosmology \todo{ctie} using the Core Cosmology Library\footnote{\todo{add url}} \todo{cite}. The value of the spin-2 shear field at each source is then calculated through an inverse spherical harmonic transform using the irregular-grid method of \todo{cite} implemented in the {\tt ducc} library\footnote{\todo{ref}}. In addition to the cosmic shear signal, we add a shape noise component, uncorrelated between sources, with a standard deviation per component of $\sigma_e=0.27$. We will quantify the validity of our approximation both with and without this component, in order to better test the contributions from the different signal-like and noise-like mode-coupling coefficients in Eq. \ref{eq:main_result}.

    The sample consists of about $2.5\times10^7$ sources across approximately $4000\,{\rm deg}^2$, and the median inter-particle distance is $\theta_{\rm ip}=0.15\,{\rm arcmin}$. Hence, power spectra can be reconstructed up to very large multipoles. To stress-test our approach, we measure power spectra up $\ell_{\rm max}=10{,}000$. Although most cosmological shear analyses typically focus on larger scales \todo{cites}, several works have shown that smaller, highly non-linear scales, may be exploited given a sufficiently accurate and flexible parametrisation for non-linearities and baryonic effects \cite{2303.05537,2403.13794}. We measure the cosmic shear power spectrum in 1000 simulations using 38 bandpowers spanning multipoles $\ell\leq\ell_{\rm max}$. We use the same binning scheme adopted in \cite{2403.13794}, with linearly-spaced bins up to $\ell=240$ and logarithmically afterwards. The simulation measurements are then used to construct the sample covariance, which we compare with our analytical estimate. The analytical covariance is calculated following Eq. \ref{eq:main_result}, using the input power spectrum to estimate $\bar{C}^{ab}_\ell$, and with the weighted variance map (Eq. \ref{eq:awa2}) calculated from one of the simulations, chosen at random.

    The result of this exercise is shown in Fig. \ref{fig:cov_shear} for both $E$-mode and $B$-mode power spectra. We find that the analytical covariance is able to describe the sample covariance at accurately on all scales in the case of $E$-modes, and on scales $\ell\gtrsim10$ for $B$-modes. The off-diagonal elements of the covariance are small (a few percent of the main diagonal), and are well characterised by the analytical estimate. Remarkably, these results hold also for noiseless simulations, where the contribution from the noise-like components of the covariance (which are exact in the analytical approach) are suppressed.

  \subsection{Galaxy clustering}\label{ssec:val.clust}
    \begin{figure}
      \centering
      \includegraphics[width=0.5\textwidth]{figures/cov_lrgs.pdf}
      \caption{Power spectrum covariance for the LRG-like simulated galaxy clustering sample. Results are shown for power spectra calculated using random catalogues to quantify the sample completeness (red) and using a continuous mask (blue). For clarity, the mask-based results have been rescaled by a factor 1.5. The points show the sample variance estimated from 1000 realisations, with the analytical Gaussian covariance as solid lines. The dashed line and crosses show the first off-diagonal element of the covariance matrix.}
      \label{fig:cov_lrgs}
    \end{figure}
    In our final example, we explore the case of galaxy clustering in a realistic setting. As discussed in Section \todo{ref}, this is an important case to test separately, since the noise-like contributions to the covariance are driven by the Poisson-like shot noise of the sample, and not by the variance of any sampled field.

    As a realistic setup, we generate mock realisations of a galaxy clustering sample with properties (number density, clustering bias) of the photometric luminous red galaxy (LRG) sample of the Legacy Survey \todo{cites} used in the analysis of \todo{cite}. To generate mock realisations, we follow the steps outlined below:
    \begin{enumerate}
      \item We start by calculating the angular power spectrum of the sample. As before, we use CCL, assuming the redshift distribution of the second redshift bin in \todo{cite}, and a galaxy bias of $b=2$. We calculate this power spectrum up to $\ell_{\rm max}=1{,}000$.
      \item We generate a Gaussian realisation of the galaxy overdensity $\delta_G$ from this power spectrum on a HEALPix map with resolution parameter $N_{\rm side}=512$.
      \item In order to generate a positive-definite overdensity map $\delta_g$ from which we can generate a Poisson-sampled catalogue, we perform a log-normal transformation \todo{cite} of the form
      \begin{equation}
        1+\delta_g=\exp[\delta_G-\sigma_G^2/2],
      \end{equation}
      where $\sigma_G^2\equiv\langle\delta_G^2\rangle$ is the variance of $\delta_G$.
      \item We generate a simulated catalog via Poisson-sampling, where the probablity of finding a galaxy in pixel $p$ is
      \begin{equation}
        p_p\propto w_p\,(1+\delta_{g,p}),
      \end{equation}
      where $w_p$ is a completeness map. The completeness map $w_p$ was reconstructed from the publicly available random catalogues used in \todo{cite}.
    \end{enumerate}
    As in previous sections, we generate 1000 such simulations and estimate the power spectrum of each of them to calculate the sample covariance. We calculate these power spectra using two different approaches: 
    A few subtlet
 

\section{Conclusion}\label{sec:conc}
    \lipsum[11]

\section*{Acknowldedgements}
  We
\emph{Software}:  We made extensive use of the {\tt numpy} \citep{oliphant2006guide, van2011numpy}, {\tt scipy} \citep{2020SciPy-NMeth}, {\tt astropy} \citep{1307.6212, 1801.02634}, {\tt healpy} \citep{Zonca2019}, {\tt matplotlib} \citep{2007CSE.....9...90H} and {\tt GetDist} \citep{Lewis:2019xzd} python packages, as well as the {\tt HEALPix} package \cite{astro-ph/0409513}.

This paper makes use of software developed for the Large Synoptic Survey Telescope. We thank the LSST Project for making their code available as free software at \url{http://dm.lsst.org}. 

\emph{Data:} 

%==============%
%   APPENDIX   %
%==============%
\onecolumngrid
\newpage

\appendix
\section{Spin-$2$ fields}\label{app:spin_s}
  For simplicity, our discussion in Sections \ref{sec:pcls} and \ref{sec:theo} has focused on spin-0 fields. The generalisation of the results presented here to the case of spin-2 fields offers no major difficulty, beyond those already described in the literature. The modifications to the coupling coefficients appearing in Eq. \ref{eq:main_result} to accommodate different spins are described in \cite{1906.11765}, with the extensions to the noise-like contributions introduced in \cite{2010.09717}. The shot noise contribution to the field's pseudo-$C_\ell$, in the case of auto-spectra, as described in \cite{2407.21013}, is given by
  \begin{equation}
    \tilde{N}^{\alpha\beta}_a=\frac{\delta^{\alpha\beta}}{4\pi}\sum_iw_i^2\frac{a_{1,i}^2+a_{2,i}^2}{2},
  \end{equation}
  where $\alpha$ and $\beta$ label the different $E$ and $B$ components of the harmonic coefficients, and $a_{1/2,i}$ are the two different components of the field at the $i$-th source.

  The general covariance of spin-2 pseudo-$C_\ell$s must account for the polarisations ($E$ or $B$) of the different fields being correlated. The calculation is relatively straightforward, although it requires non-trivial amounts of bookkeeping. We reproduce the main results here for completeness.  As in the case of scalar fields, the first step is the calculation of a general two-point correlator of masked spin-2 fields (Eq. \ref{eq:general_2pt_cat}):
  \begin{equation}
    \langle\tilde{a}^\alpha_{\ell m}\tilde{a}^\beta_{\ell'm'}\rangle=\sum_{i\neq j}(\cdots)+\,_2{\cal M}^{\alpha\beta}_{\ell m,\ell'm'}(\overline{(aw_a)^2}),
  \end{equation}
  where the spin-2 map-level coupling coefficients are defined similarly to the spin-0 ones (Eq. \ref{eq:map_mcm}):
  \begin{equation}
    _2{\cal M}^{\alpha\beta}_{\ell m,\ell'm'}(w)\equiv\int d\nv\,w(\nv)\,\left(_2{\sf Y}^\dagger_{\ell m}(\nv)\,\,_2{\sf Y}_{\ell'm'}(\nv)\right)_{\alpha\beta},
  \end{equation}
  and $_2{\sf Y}_{\ell m}(\nv)$ are the spin-2 spherical harmonic matrices defined in Appendix A of \cite{1809.09603}. With this in hand, and following the same approximations outlined in \cite{1906.11765}, we can calculate the final covariance matrix, including contributions from both self-pairs and different pairs. Concretely, the spin-$2$ version of Eq. \ref{eq:main_result} reads
  \begin{align}
    {\rm Cov}(\tilde{C}^{a_\alpha b_\beta}_\ell,\tilde{C}^{f_\phi g_\gamma}_{\ell'})=&\,\overline{C}^{a_{\alpha'}f_{\phi'}}_{(\ell,\ell')}\overline{C}^{b_{\beta'}g_{\gamma'}}_{(\ell,\ell')}\left[\Xi^{\alpha\phi,\alpha'\phi'}_{\beta\gamma,\beta'\gamma'}\right]_{\ell\ell'}\left(\bar{w}_a\bar{w}_f-\delta^K_{af}\overline{w_a^2},\bar{w}_b\bar{w}_g-\delta^K_{bg}\overline{w_b^2}\right)\\
    &+\delta^K_{bg}\,\overline{C}^{a_{\alpha'}f_{\phi'}}_{(\ell,\ell')}\left[\Xi^{\alpha\phi,\alpha'\phi'}_{\beta\gamma}\right]_{\ell\ell'}\left(\bar{w}_a\bar{w}_f-\delta^K_{af}\overline{w_a^2},\overline{(bw_b)^2}\right)\\
    &+\delta^K_{af}\,\overline{C}^{b_{\beta'}g_{\gamma'}}_{(\ell,\ell')}\left[\Xi^{\alpha\phi}_{\beta\gamma,\beta'\gamma'}\right]_{\ell\ell'}\left(\overline{(aw_a)^2},\bar{w}_b\bar{w}_g-\delta^K_{bg}\overline{w_b^2}\right)\\
    &+\delta^K_{af}\,\delta^K_{bg}\left[\Xi^{\alpha\phi}_{\beta\gamma}\right]_{\ell\ell'}\left(\overline{(aw_a)^2},\overline{(bw_b)^2}\right)\\\nonumber
      &+(c\leftrightarrow d).
  \end{align}
  Here, all repeated indices are implicitly summed over, and the spin-2 covariance coupling terms for the signal-noise ($\Xi^{\alpha\phi,\alpha'\phi'}_{\beta\gamma}$), noise-signal ($\Xi^{\alpha\phi}_{\beta\gamma,\beta'\gamma'}$), and noise-noise ($\Xi^{\alpha\phi}_{\beta\gamma}$) contributions are given by:
  \begin{align}
    &\left[\Xi^{EE}_{EE}\right]_{\ell\ell'}=
    \left[\Xi^{EE}_{BB}\right]_{\ell\ell'}=
    \left[\Xi^{BB}_{EE}\right]_{\ell\ell'}=
    \left[\Xi^{BB}_{BB}\right]_{\ell\ell'}=\,_2\Xi_{\ell\ell'}^+,\\
    &\left[\Xi^{EB}_{EB}\right]_{\ell\ell'}=
    \left[\Xi^{BE}_{BE}\right]_{\ell\ell'}=
    -\left[\Xi^{EB}_{BE}\right]_{\ell\ell'}=
    -\left[\Xi^{BE}_{EB}\right]_{\ell\ell'}=\,_2\Xi_{\ell\ell'}^-,\\
    &\left[\Xi^{EE}_{XY,XY}\right]_{\ell\ell'}=\left[\Xi^{BB}_{XY,XY}\right]_{\ell\ell'}=\,_2\Xi^+_{\ell\ell'},\\
    &\left[\Xi^{EB}_{EE,EB}\right]_{\ell\ell'}=\left[\Xi^{EB}_{BB,EB}\right]_{\ell\ell'}=\left[\Xi^{EB}_{BE,EE}\right]_{\ell\ell'}=\left[\Xi^{EB}_{BE,BB}\right]_{\ell\ell'}=\,_2\Xi^-_{\ell\ell'},\\
    &\left[\Xi^{BE}_{\beta\gamma,\beta'\gamma'}\right]_{\ell\ell'}=-\left[\Xi^{EB}_{\beta\gamma,\beta'\gamma'}\right]_{\ell\ell'}=\left[\Xi^{BE}_{\beta'\gamma',\beta\gamma}\right]_{\ell\ell'},\hspace{12pt}\left[\Xi^{\beta\gamma,\beta'\gamma'}_{\alpha\phi}\right]_{\ell\ell'}=\left[\Xi^{\alpha\phi}_{\beta\gamma,\beta'\gamma'}\right]_{\ell\ell'},
  \end{align}
  and the signal-signal terms ($\Xi^{\alpha\phi,\alpha'\phi'}_{\beta\gamma,\beta'\gamma'}$) are provided in Eqs. 2.63--2.65 of \cite{1906.11765}. Here, $_2\Xi^\pm_{\ell\ell'}$ are the spin-2 mode-coupling coefficients:
  \begin{equation}
    _2\Xi^{\pm}_{\ell\ell'}(v,w)=\sum_{\ell''}\frac{2\ell''+1}{4\pi}\tilde{C}^{v,w}_{\ell''}\frac{1\pm(-1)^{L}}{2}\wtj{\ell}{\ell'}{\ell''}{2}{-2}{0}^2,
  \end{equation}
  where $L\equiv\ell+\ell'+\ell''$.

  The only remaining piece is the procedure used to calculate the local weighted variance map, $\overline{(aw_a)^2}$ (see Eq. \ref{eq:self_var}), used to calculate the coupling coefficients of the noise-like contributions to the covariance. As we discussed at the end of Section \ref{ssec:theo.general}, in the case of spin-0 fields, the only possible estimate of the local variance is $\langle a_i^2\rangle\simeq a_i^2$, in the case of spin-2 fields, we can follow from the definition of $\tilde{N}_a$ above and use $\langle a_i^2\rangle\simeq(a_{1,i}^2+a_{2,i}^2)/2$ as a local variance estimate. I.e. in the case of spin-2 fields, Eq. \ref{eq:awa2} reads
  \begin{equation}
      \left(\overline{(aw_a)^2}\right)_{\ell m}\equiv\sum_iw_i^2\frac{a_{1,i}^2+a_{2,i}^2}{2}Y_{\ell m,i}^*.
  \end{equation}
  With this choice, the auto-spectrum of $\overline{(aw_a)^2}$ must be corrected by the following additive constant:
  \begin{equation}
    \tilde{N}_{\overline{(aw_a)^2}}=\frac{1}{4\pi}\sum_iw_i^4\left(\frac{a_{1,i}^2+a_{2,i}^2}{2}\right)^2.
  \end{equation}
  This generalises Eq. \ref{eq:4point_corr_spin0} for non-scalar fields.

\section{De-biasing $\tilde{C}_\ell^{\overline{(aw_a)^2}}$}\label{app:awa2_shotnoise}
  We derive here the bias to the pseudo-$C_\ell$ of the local variance map $\overline{(aw_a)^2}$, defined in Eq. \ref{eq:self_var}, when the local variance $\sigma_i^2\equiv \langle a_i^2\rangle$ is estimated directly from the data, and accounting for the subtraction of $\tilde{N}_a$ when estimating the auto-spectra of catalog-based fields.

  \subsection{Correcting for $\langle a_i^2\rangle\neq a_i^2$}
    The local variance field $\overline{(aw_a)^2}$, entering the noise-noise component of the covariance matrix (Eq. \ref{eq:main_result}) through its pseudo-$C_\ell$ $\tilde{C}_\ell^{\overline{(aw_a)^2}}$, is defined in  Eq. \ref{eq:self_var} in terms of the point-wise variance $\langle a_i^2\rangle$. For a given catalog, assuming that we do not have an a-priori estimate of $\langle a_i^2\rangle$, we replace it simply by its point estimate $a_i^2$. The data-based local variance map is then given by Eq. \ref{eq:awa2}. The ensemble average of the pseudo-$C_\ell$ of this map is:
    \begin{align}\nonumber
      \left\langle\hat{\tilde{C}}^{\overline{(aw_a)^2}}_\ell\right\rangle
      &= \sum_{ij}w_i^2w_j^2\langle a_i^2a_j^2\rangle \frac{1}{2\ell+1}\sum_mY^*_{\ell m,i}Y_{\ell m,j}\\\nonumber
      &= \sum_{ij}w_i^2w_j^2\langle a_i^2\rangle\langle a_j^2\rangle \frac{1}{2\ell+1}\sum_mY^*_{\ell m,i}Y_{\ell m,j}+2\sum_iw_i^4\langle a_i^2\rangle^2\,\frac{1}{2\ell+1}\sum_m|Y^*_{\ell m,i}|^2\\\label{eq:pcl_var_corr1}
      &=\tilde{C}^{\overline{(aw_a)^2}}_\ell+2\frac{1}{4\pi}\sum_iw_i^4\langle a_i^2\rangle^2
    \end{align}
    where we have expanded the four-point cumulant in the second line using Wick's theorem, assuming uncorrelated field values ($\langle a_ia_j\rangle=\delta_{ij}\langle a_i^2\rangle$). The pseudo-$C_\ell$ of the data-based local variance map is thus biased by the second term in the equation above. We must therefore correct for this term, which we can do, again, through a data-based estimate of the field's squared variance $\langle a^2_i\rangle^2$. Since, for Gaussian fields, $\langle a_i^4\rangle=3\langle a_i^2\rangle^2$, this data-based correction can be calculated as:
    \begin{equation}
      \Delta\tilde{C}^{\overline{(aw_a)^2},1}_\ell=\frac{2}{3}\frac{1}{4\pi}\sum_iw_i^4a_i^4
    \end{equation}

  \subsection{Correcting for $\tilde{N}_a$}
    Our derivation of the Gaussian iNKA covariance in Section \ref{ssec:theo.general} assumed, for simplicity, that we are interested in the covariance of the pseudo-$C_\ell$ $\tilde{C}^{ab}_\ell$. While in general this is true, as we described in Section \ref{ssec:pcls.cat}, the auto-correlation of catalog-based fields must be corrected for the contribution from self-pairs $\tilde{N}_a$ (Eq. \ref{eq:Na}). For auto-correlations, the quantity of interest is thus the combination
    \begin{equation}
      \tilde{C}^{aa,{\rm corr}}_\ell\equiv \tilde{C}^{aa}_\ell-\tilde{N}_a=\frac{1}{2\ell+1}\sum_m|\tilde{a}_{\ell m}|^2-\frac{1}{4\pi}\sum_i(w_ia_i)^2.
    \end{equation}
    The covariance of this quantity is then:
    \begin{align}\nonumber
      {\rm Cov}(\tilde{C}^{aa,{\rm corr}}_\ell,\tilde{C}^{aa,{\rm corr}}_{\ell'})=\frac{1}{(2\ell+1)(2\ell'+1)}&\left[\sum_{mm'}\langle |\tilde{a}_{\ell m}|^2|\tilde{a}_{\ell' m'}|^2\rangle-\frac{2\ell+1}{4\pi}\sum_{m'}\sum_iw_i^2\langle a_i^2|\tilde{a}_{\ell' m'}|^2\rangle\right.\\\nonumber
      &\left.-\frac{2\ell'+1}{4\pi}\sum_m\sum_iw_i^2\langle a_i^2|\tilde{a}_{\ell m}|^2\rangle+\frac{(2\ell+1)(2\ell'+1)}{(4\pi)^2}\sum_{ij}w_i^2w_j^2\langle a_i^2a_j^2\rangle
      \right].
    \end{align}
    Using Wick's theorem on all the 4-point correlators above for a Gaussian, noise-like field ($\langle a_ia_j\rangle=\delta_{ij}\langle a_i^2\rangle$), the expression above can be simplified, after a bit of algebra, into:
    \begin{align}
      {\rm Cov}(\tilde{C}^{aa,{\rm corr}}_\ell,\tilde{C}^{aa,{\rm corr}}_{\ell'})
      &=2\left[\Xi_{\ell\ell'}\left(\overline{(aw_a)^2},\overline{(aw_a)^2}\right)-\frac{1}{(4\pi)^2}\sum_i w_i^4\langle a_i^2\rangle^2\right].
    \end{align}
    This is equivalent to the noise-noise contribution in Eq. \ref{eq:main_result}, corrected by the second term above. Further insight into this correction may be gained by expanding the coupling coefficient $\Xi_{\ell\ell'}$ in terms of the pseudo-$C_\ell$ of its arguments (see Eq. \ref{eq:Xi}):
    \begin{align}\nonumber
      {\rm Cov}(\tilde{C}^{aa,{\rm corr}}_\ell,\tilde{C}^{aa,{\rm corr}}_{\ell'})
      &=2\left[\sum_{\ell''}\frac{2\ell''+1}{4\pi}\wtj{\ell}{\ell'}{\ell''}{0}{0}{0}^2\,\tilde{C}^{\overline{(aw_a)^2}}_{\ell''}-\frac{1}{(4\pi)^2}\sum_i w_i^4\langle a_i^2\rangle^2\right]\\
      &=2\sum_{\ell''}\frac{2\ell''+1}{4\pi}\wtj{\ell}{\ell'}{\ell''}{0}{0}{0}^2\left[\tilde{C}^{\overline{(aw_a)^2}}_{\ell''}-\frac{1}{4\pi}\sum_i w_i^4\langle a_i^2\rangle^2\right],
    \end{align}
    where, in the second line, we have used the property
    \begin{equation}
      \sum_{\ell''}(2\ell''+1)\wtj{\ell}{\ell'}{\ell''}{0}{0}{0}^2=1.
    \end{equation}

    We thus see that, to account for the $\tilde{N}_a$ subtraction in auto-spectra, the power spectrum of the local variance field entering the noise-noise covariance must be corrected by a term that is half of the correction derived in the previous section, which accounts for the fact that the local variance map is estimated from the data. Note that this correction is exactly the contribution to the pseudo-$C_\ell$ of $\overline{(aw_a)^2}$ from self-pairs. I.e. it plays the same role as $\tilde{N}_a$ and $\tilde{N}_w$, accounting for the ``shot noise'' contribution to their corresponding power spectra. As in the previous section, we can estimate this term from the data, obtaining the second correction
    \begin{equation}
      \Delta\tilde{C}^{\overline{(aw_a)^2},2}_\ell=\frac{1}{3}\frac{1}{4\pi}\sum_iw_i^4a_i^4.
    \end{equation}

    Thus we find that, when calculating the noise-noise contribution to the covariance matrix of auto-correlations, the pseudo-$C_\ell$ of the local variance map must be corrected by subtracting the quantity
    \begin{equation}
      \Delta\tilde{C}^{\overline{(aw_a)^2},1}_\ell+\Delta\tilde{C}^{\overline{(aw_a)^2},2}_\ell=\frac{1}{4\pi}\sum_iw_i^4a_i^4\equiv\tilde{N}_{\overline{(aw_a)^2}},
    \end{equation}
    defined in Eq. \ref{eq:4point_corr_spin0}.

% \section{$\chi^2$ validation for iNKA covariances}\label{app:inka_chi2}

% \begin{figure}
%   \centering
%   \includegraphics[width=\textwidth]{figures/chi2_spin0_alldiff_paper.png}
%   \caption{$\chi^2$ distributions for scalar field validation cases, with none of the catalogs sharing any sources.}\label{fig:chi2_spin0_alldiff}
% \end{figure}
% \begin{figure}
%   \centering
%   \includegraphics[width=\textwidth]{figures/chi2_spin0_allsame_paper.png}
%   \caption{$\chi^2$ distributions for scalar field validation cases, with all catalogs sharing the same sources.}\label{fig:chi2_spin0_allsame}
% \end{figure}
% \begin{figure}
%   \centering
%   \includegraphics[width=\textwidth]{figures/chi2_spin2_alldiff_paper.png}
%   \caption{$\chi^2$ distributions for spin-2 field validation cases, with none of the catalogs sharing any sources.}\label{fig:chi2_spin0_alldiff}
% \end{figure}
% \begin{figure}
%   \centering
%   \includegraphics[width=\textwidth]{figures/chi2_spin2_allsame_paper.png}
%   \caption{$\chi^2$ distributions for spin-2 field validation cases, with all catalogs sharing the same sources.}\label{fig:chi2_spin2_allsame}
% \end{figure}
% \begin{figure}
%   \centering
%   \includegraphics[width=0.49\textwidth]{figures/chi2_spin2xspin0_alldiff_paper.png}
%   \includegraphics[width=0.49\textwidth]{figures/chi2_spin2xspin0_allsame_paper.png}
%   \caption{$\chi^2$ distributions for spin-2 $\times$ spin-0 field validation cases. {\it} Left:} Case where none of the catalogs share the same source positions. {\it Right:} Case where all catalogs share the same source positions. \label{fig:chi2_spin2xspin0_allsame}
% \end{figure}



\section{Numerical implementation of the brute force calculation} \label{app:brute_code}
In this section, we explain the code implementation and optimisation to compute the covariance matrix. We implement the covariance calculation using the \texttt{jax} library and design the code to run optimally on a GPU. The main advantage is to exploit multiple properties that \texttt{jax} offers to accelerate the covariance calculation without approximation.

The real bottleneck of this calculation is calculating the highly dimensional $P_{\ell, ij}$, which have dimensions $\ell_{max} \times N_{\text{FRBs}} \times N_{\text{FRBs}}$ in addition to the expensive calculation of the Legendre polynomials for large $\ell$. For the former problem, we leverage the \texttt{jax} package to accelerate matrix computations. For the latter, we use the recursion relation of the Legendre Polynomials. We can also accelerate the calculation of the sum over Legendre polynomials using the Bonnet recurrence formula  to produce recursive calculations of the Legendre polynomial \citep{Hademenos_Spiegel_Liu_2001}
\begin{equation}
    P_{n}(x) = \frac{2n-1}{n} \, x P_n(x) - \frac{n - 1}{n} P_{n-2}(x).
\end{equation}
For consecutive calculation, we use \texttt{jax.lax} property to loop over the calculation for $\ell_\mathrm{min} \leq \ell \leq \ell_\mathrm{max}$. In addition, to calculate the correlator brackets given in Eq. \ref{eq:cov_brute}, we use the binned estimated spectra $S^B_\ell$ by assuming that the value of the power spectrum is constant for different multipoles for a particular bandpower $B$, i.e.
\begin{equation}
    \sum_{\ell} \frac{2\ell+1}{4\pi} S_\ell P_{\ell,ij}  \approx \sum_{B} S^B_\ell  \sum_{\ell \in B} \frac{2\ell+1}{4\pi}  P_{\ell,ij}.
\end{equation}

When calculating  a binned covariance matrix, we introduce a binning matrix based on the given bandpower edges and apply it to the Legendre polynomials $P_{\ell, ij}$ before summing over the sources, such as to reduce the number of operations. In other words, the binned covariance matrix is given as

\begin{equation}
    \mathrm{Cov}(\hat{C}_\ell^{aa}, \hat{C}_{\ell'}^{aa})  =  2
    \sum_{i \not= j,\, p \not= q}\frac{w_i w_jw_p w_q}{(4\pi)^2}   [BM^{-1}P]_{\ell, ij} [PM^{-1}B]_{\ell', pq}
     \left\langle a_i a_p \right\rangle \left\langle a_j a_q \right\rangle.
\end{equation}

Depending on the memory restrictions of the available GPU and the number of sources, it is possible to accommodate different memory restrictions by splitting the sum into smaller parts. 

The code to perform brute-force covariance matrix calculations is made available at this URL \todo{add URL}. The repository contains two implementations of the covariance matrix: one that can be used in a for-loop for a situation with small GPU memory or a large number of sources, and another for a small number of sources that can be processes with the memory of a single GPU. 

\section{Reduction of the direct summation covariance to the NKA limit}\label{app:brute_nka}
The brute-force expression, Eq.~\ref{eq:cov_brute}, is exact for Gaussian fields, whereas the analytical covariance of Sect.~\ref{sec:theo} relies on the NKA. Here we show that the former reduces to the latter in the appropriate limit.
Using the addition theorem, the direct sum can be expressed through the masked harmonic coefficients $\tilde{a}_{\ell m}$ as:
\begin{equation}\label{eq:brute_as_cov}
    \mathrm{Cov}(\hat{C}_\ell^{aa},\hat{C}_{\ell'}^{aa}) = \sum_{\ell_1,\ell_2}[M^{-1}]_{\ell\ell_1}[M^{-1}]_{\ell_2\ell'}\,\mathrm{Cov}(\tilde{C}_{\ell_1}^{aa},\tilde{C}_{\ell_2}^{aa}),
\end{equation}
where
\begin{equation}\label{eq:pcl_cov_harm}
    \mathrm{Cov}(\tilde{C}_{\ell_1}^{aa},\tilde{C}_{\ell_2}^{aa}) = \frac{2}{(2\ell_1+1)(2\ell_2+1)}\sum_{m_1 m_2}\left|\langle \tilde{a}_{\ell_1 m_1}\tilde{a}^*_{\ell_2 m_2}\rangle\right|^2.
\end{equation}
This is the Gaussian pseudo-$C_\ell$ covariance of Sect.~\ref{ssec:pcls.cov}. The brute-force result is therefore the decoupled pseudo-$C_\ell$ covariance evaluated without approximation. The restriction to distinct source pairs ($i\neq j$, $p\neq q$) corresponds to the self-pair debiasing already applied in Eq.~\ref{eq:invert}.
 
Inserting the field two-point function, Eq.~\ref{eq:field_corr}, into the spherical harmonic expansion of the masked field separates the distinct-source signal from the self-pair shot noise. As in Sect.~\ref{sec:theo}, the distinct ($i\neq j$) pairs carry only the signal and, under the NKA, produce the self-pair-subtracted mask $\bar{w}_a^2-\overline{w_a^2}$ together with the iNKA spectrum $\overline{C}^{aa}$ of Eq.~\ref{eq:clbar0}, while the self-pairs ($i=j$) carry the full local variance $\langle a_i^2\rangle$ and assemble into the shot-noise map $\overline{(aw_a)^2}$ of Eq.~\ref{eq:awa2},
\begin{equation}\label{eq:masked_corr}
    \langle\tilde{a}_{\ell_1 m_1}\tilde{a}^*_{\ell_2 m_2}\rangle \approx \overline{C}^{aa}_{(\ell_1,\ell_2)}\,{\cal M}_{\ell_1 m_1,\ell_2 m_2}\!\left(\bar{w}_a^2-\overline{w_a^2}\right) + {\cal M}_{\ell_1 m_1,\ell_2 m_2}\!\left(\overline{(aw_a)^2}\right).
\end{equation}
Here we have pulled the slowly-varying $C^{aa}_{\ell''}\to\overline{C}^{aa}_{(\ell_1,\ell_2)}$ out of the convolution. Squaring Eq.~\ref{eq:masked_corr} and using
\begin{equation}\label{eq:M2_xi}
    \sum_{m_1 m_2}{\cal M}_{\ell_1 m_1,\ell_2 m_2}(u)\,{\cal M}^*_{\ell_1 m_1,\ell_2 m_2}(v) = (2\ell_1+1)(2\ell_2+1)\,\Xi_{\ell_1\ell_2}(u,v),
\end{equation}
Eq.~\ref{eq:pcl_cov_harm} becomes
\begin{equation}\label{eq:brute_nka}
\begin{split}
    \mathrm{Cov}(\tilde{C}_{\ell}^{aa},\tilde{C}_{\ell'}^{aa}) \approx 2\Big[&\big(\overline{C}^{aa}_{(\ell,\ell')}\big)^2\,\Xi_{\ell\ell'}\!\left(\bar{w}_a^2-\overline{w_a^2},\,\bar{w}_a^2-\overline{w_a^2}\right) + 2\,\overline{C}^{aa}_{(\ell,\ell')}\,\Xi_{\ell\ell'}\!\left(\bar{w}_a^2-\overline{w_a^2},\,\overline{(aw_a)^2}\right) \\
    &+ \Xi_{\ell\ell'}\!\left(\overline{(aw_a)^2},\,\overline{(aw_a)^2}\right)\Big],
\end{split}
\end{equation}
which is Eq.~\ref{eq:main_result} for $a=b=c=d$. 
 



\twocolumngrid

\bibliography{main}

\end{document}


